\documentclass[11pt]{article}
\usepackage[utf8]{inputenc}
\usepackage[T1]{fontenc}
\usepackage{lmodern}
\usepackage[margin=1in]{geometry}
\usepackage{amsmath,amssymb}
\usepackage{enumitem}

\title{}
\author{}
\date{}

\begin{document}
\maketitle

\section*{Question 1: Write each set in roster notation}

\begin{enumerate}
    \item[(a)] The set of all integers less than \(7\) is infinite and extends without limit in the negative direction. In roster notation:
    \[
    \{\ldots,-4,-3,-2,-1,0,1,2,3,4,5,6\}.
    \]
    \item[(b)] The odd integers from \(1\) to \(17\), inclusive, are:
    \[
    \{1,3,5,7,9,11,13,15,17\}.
    \]
\end{enumerate}

\section*{Question 2: Set operations}

Given:
\[
A=\{1,2,3,4\},\qquad
B=\{2,4,6,8\},\qquad
C=\{3,4,5,6\}.
\]

\begin{enumerate}
    \item[(a)] Find \(B\cup C\). The union contains every element in either set, without repetition:
    \[
    B\cup C=\{2,3,4,5,6,8\}.
    \]
    \item[(b)] Find \(A\cap(B-C)\). First,
    \[
    B-C=\{2,8\},
    \]
    because \(4\) and \(6\) are also in \(C\). Then
    \[
    A\cap\{2,8\}=\{2\}.
    \]
\end{enumerate}

\section*{Question 3: Find the non-zero pre-image of \(4\)}

Given \(f(x)=x^2+2x+4\). Set \(f(x)\) equal to \(4\):
\begin{align*}
x^2+2x+4 &= 4,\\
x^2+2x &= 0,\\
x(x+2) &= 0.
\end{align*}
Therefore, \(x=0\) or \(x=-2\). The non-zero pre-image is \(-2\).

\[
\boxed{x=-2}
\]

\section*{Question 4: Students playing football and hockey}

Let \(F\) be the set of students who play football and \(H\) the set who play hockey. The number who play at least one sport is
\[
n(F\cup H)=n(F)+n(H)-n(F\cap H)=19+13-8=24.
\]
Add the \(7\) students who play neither sport:
\[
24+7=31.
\]

\[
\boxed{31\text{ students}}
\]

\section*{Question 5: Carlsberg Malawi company survey}

The number who liked green or brown is
\[
220+270-70=420.
\]
Add the \(50\) people who liked neither type:
\[
420+50=470
\]
people interviewed. The total cost is
\[
470\times \mathrm{K}7{,}655=\mathrm{K}3{,}597{,}850.
\]

\[
\boxed{\mathrm{K}3{,}597{,}850}
\]

\section*{Question 6: Medals in three categories}

The total number of category memberships is
\[
36+12+18=66.
\]
Let \(x\) be the number of people receiving medals in exactly one category, and \(y\) the number receiving medals in exactly two categories. Four people received medals in all three categories. Since \(45\) people received medals,
\[
x+y+4=45,
\qquad\text{so}\qquad
x+y=41.
\]
Counting category memberships gives
\[
x+2y+3(4)=66,
\qquad\text{so}\qquad
x+2y=54.
\]
Subtracting the first equation from the second gives \(y=13\).

\[
\boxed{13\text{ people received medals in exactly two categories}}
\]

\section*{Question 7: Solve the equation}

Given \(f(x)= 4x-7\) and \(g(x)=x+2\), solve \(f(2x)=7g(x)\).

\textit{Working}
\[
f(x) = 4x - 7
\]
\[
g(x) = x + 2
\]

\textbf{Step 1: Find \(f(2x)\)}

\[
\begin{aligned}
f(2x) &= 4(2x) - 7 \\
      &= 8x - 7
\end{aligned}
\]

\textbf{Step 2: Find \(7g(x)\)}

\[
\begin{aligned}
7g(x) &= 7(x+2) \\
     &= 7x + 14
\end{aligned}
\]

\textbf{Step 3: Equate the expressions}

\[
8x - 7 = 7x + 14
\]

\[
x = 21
\]

\[
\boxed{x = 21}
\]

\section*{Question 8: Conjunction (AND)}

The conjunction \(p\land q\) is true only when both \(p\) and \(q\) are true. Here, \(p\) means ``John scored a goal'' and \(q\) means ``John won the game.''

\[
\begin{array}{c|c|c}
p & q & p\land q\\
\hline
T & T & T\\
T & F & F\\
F & T & F\\
F & F & F
\end{array}
\]

\section*{Question 9: Birthday-month deposits}

The deposits are
\[
\mathrm{K}1+\mathrm{K}2+\mathrm{K}3+\cdots+\mathrm{K}30.
\]
Using the sum formula for the first \(n\) positive integers,
\[
S=\frac{n(n+1)}{2}.
\]
Therefore,
\[
S=\frac{30(31)}{2}=465.
\]

\[
\boxed{\mathrm{K}465}
\]

\section*{Question 10: Find the 200th term of an arithmetic sequence}

The fifth term is \(23\) and the sixth term is \(25\), so the common difference is
\[
d=25-23=2.
\]
Using \(T_n=a+(n-1)d\), the fifth term gives
\[
a+4(2)=23,
\qquad\text{so}\qquad
a=15.
\]
Thus,
\[
T_{200}=15+199(2)=413.
\]

\[
\boxed{\text{The 200th term is }413}
\]

\section*{Question 11: Find the eighth term of an arithmetic sequence}

The fourth term is \(7\) and the fifth term is \(4\), so
\[
d=4-7=-3.
\]
Using \(T_n=a+(n-1)d\), the fourth term gives
\[
a+3(-3)=7,
\qquad\text{so}\qquad
a=16.
\]
Therefore,
\[
T_8=16+7(-3)=-5.
\]

\[
\boxed{\text{The eighth term is }-5}
\]

\section*{Question 12: Substrate-level phosphorylation and chemiosmosis}

\paragraph{Substrate-level phosphorylation}
This is the direct formation of ATP when an enzyme transfers a phosphate group from a high-energy substrate to ADP. It occurs in glycolysis in the cytoplasm and in one reaction of the citric acid cycle in the mitochondrial matrix. It does not require ATP synthase or a proton gradient.

\paragraph{Chemiosmosis in oxidative phosphorylation}
Electrons move through the electron transport chain in the inner mitochondrial membrane. The released energy pumps protons across the membrane, creating an electrochemical gradient. Protons flow back through ATP synthase, which uses this energy to produce ATP. Oxygen is the final electron acceptor in the electron transport chain.

\paragraph{Comparison}
Both processes produce ATP. Substrate-level phosphorylation transfers phosphate directly to ADP, whereas chemiosmosis uses the energy of a proton gradient and requires ATP synthase. Chemiosmosis is part of oxidative phosphorylation and depends indirectly on oxygen.

\section*{Question 13: Why the citric acid cycle is considered an aerobic pathway}

The citric acid cycle does not use oxygen directly, but it depends indirectly on oxygen to continue operating.

\begin{enumerate}
    \item \textbf{Reduced coenzymes:} The cycle produces \(\mathrm{NADH}\) and \(\mathrm{FADH_2}\), which carry high-energy electrons to the electron transport chain.
    \item \textbf{Oxygen as final electron acceptor:} Oxygen accepts electrons at the end of the electron transport chain and combines with protons to form water. Without oxygen, electron transport stops.
    \item \textbf{Regeneration of coenzymes:} The electron transport chain regenerates \(\mathrm{NAD^+}\) and \(\mathrm{FAD}\) from \(\mathrm{NADH}\) and \(\mathrm{FADH_2}\). These oxidised coenzymes are needed for the citric acid cycle to continue.
\end{enumerate}

\paragraph{Conclusion}
The citric acid cycle is therefore considered aerobic because its continued operation depends on oxygen-dependent oxidative phosphorylation.

\section*{Question 14: Osteoporosis}

\paragraph{Definition}
Osteoporosis is a skeletal disorder characterised by reduced bone mass and deterioration of bone tissue, resulting in fragile bones and an increased risk of fractures.

\subsection*{(i) Biochemical basis}

\begin{enumerate}
    \item \textbf{Imbalance between bone formation and resorption:} Bone breakdown by osteoclasts exceeds bone formation by osteoblasts.
    \item \textbf{Calcium deficiency or poor absorption:} Inadequate calcium reduces the minerals available for bone formation.
    \item \textbf{Vitamin D deficiency:} Low vitamin D reduces intestinal calcium absorption.
    \item \textbf{Hormonal changes:} Reduced oestrogen, especially after menopause, increases osteoclast activity and bone loss.
    \item \textbf{Parathyroid hormone:} Persistently increased PTH, often in response to low blood calcium, can increase bone resorption.
    \item \textbf{Reduced bone mineral density:} Loss of calcium and phosphate makes bones less dense and weaker.
    \item \textbf{Deterioration of bone structure:} Changes in the collagen framework and bone microarchitecture reduce bone strength.
\end{enumerate}

\subsection*{(ii) Clinical features}

\begin{enumerate}
    \item Back pain, particularly when vertebral compression fractures occur.
    \item Fractures after minor falls or injuries.
    \item Loss of height due to compression of vertebrae.
    \item Stooped posture or forward curvature of the upper spine (kyphosis).
    \item Hip fractures, which may occur after a fall.
    \item Wrist fractures, particularly after trying to break a fall.
    \item Reduced mobility due to pain or fractures.
    \item Difficulty performing daily activities and loss of independence.
    \item Spinal deformity resulting from repeated vertebral fractures.
    \item Some people have no symptoms until a fracture occurs.
\end{enumerate}

\paragraph{Conclusion}
Osteoporosis weakens bones because bone resorption exceeds bone formation. Adequate calcium and vitamin D, weight-bearing exercise, fall prevention, and appropriate medical care help protect bone health.

\section*{Question 15: Reactive oxygen species, superoxide ions and their harmful effects}

\subsection*{1. Introduction}
Reactive oxygen species (ROS) are chemically reactive oxygen-containing molecules formed naturally during cellular metabolism and also by exposure to certain environmental factors. In controlled amounts, they help with cell signalling and immune defence. Excessive ROS can damage cells.

\subsection*{2. Types of ROS}
\begin{itemize}
    \item \textbf{Superoxide anion} (\(\mathrm{O_2^{\bullet-}}\)): formed when oxygen gains one electron.
    \item \textbf{Hydrogen peroxide} (\(\mathrm{H_2O_2}\)): formed when superoxide is converted by superoxide dismutase. It is a ROS but not a free radical.
    \item \textbf{Hydroxyl radical} (\(\mathrm{\bullet OH}\)): extremely reactive and can damage DNA, proteins, and lipids.
    \item \textbf{Singlet oxygen} (\(\mathrm{{}^{1}O_2}\)): an excited form of oxygen that can react with biological molecules.
\end{itemize}

\subsection*{3. Superoxide ions}
Superoxide is formed when molecular oxygen gains one electron:
\[
\mathrm{O_2+e^-\longrightarrow O_2^{\bullet-}}.
\]
It can be generated in mitochondria and by enzymes involved in cellular reactions. Superoxide dismutase converts it into hydrogen peroxide and oxygen:
\[
\mathrm{2O_2^{\bullet-}+2H^+\longrightarrow H_2O_2+O_2}.
\]

\subsection*{4. Sources of ROS}
\begin{enumerate}
    \item Mitochondrial electron transport during cellular respiration.
    \item Inflammation and immune-cell activity.
    \item Ultraviolet and ionising radiation.
    \item Cigarette smoke.
    \item Air pollution and some environmental chemicals.
    \item Certain drugs and toxins.
    \item Normal metabolic oxidation reactions.
\end{enumerate}

\subsection*{5. Harmful effects on body systems}
Excessive ROS can cause oxidative stress, which occurs when ROS production exceeds antioxidant defences.

\begin{enumerate}
    \item \textbf{Cell membranes:} ROS cause lipid peroxidation, damaging membranes and changing their permeability.
    \item \textbf{DNA:} Oxidative damage can cause mutations and contribute to cancer development.
    \item \textbf{Proteins:} Oxidation can change protein structure and reduce enzyme function.
    \item \textbf{Cardiovascular system:} Oxidative stress contributes to blood-vessel dysfunction and atherosclerosis.
    \item \textbf{Nervous system:} Oxidative damage may contribute to neurodegenerative disorders such as Alzheimer's and Parkinson's diseases.
    \item \textbf{Respiratory system:} Oxidative injury contributes to inflammation and lung damage, including in chronic obstructive pulmonary disease.
    \item \textbf{Immune system:} Excessive ROS can injure immune cells and disturb normal immune responses.
    \item \textbf{Liver:} Oxidative stress can contribute to liver-cell injury and inflammation.
    \item \textbf{Kidneys:} Oxidative injury can damage kidney tissues and impair their function.
    \item \textbf{Ageing:} Accumulated oxidative damage contributes to cellular ageing.
    \item \textbf{Diabetes:} Oxidative stress contributes to blood-vessel and tissue complications associated with diabetes.
    \item \textbf{Eyes:} Oxidative damage is involved in cataract formation and other age-related eye conditions.
    \item \textbf{Inflammation:} ROS can activate inflammatory pathways and worsen tissue injury.
    \item \textbf{Cancer:} ROS-related DNA damage can contribute to mutations involved in cancer.
    \item \textbf{Muscles:} Excessive oxidative damage can affect muscle proteins and membranes, impairing muscle function.
\end{enumerate}

\subsection*{6. Antioxidant defence mechanisms}
The body uses antioxidant systems to limit ROS damage.

\paragraph{Enzymatic antioxidants}
\begin{itemize}
    \item Superoxide dismutase (SOD) converts superoxide into hydrogen peroxide and oxygen.
    \item Catalase converts hydrogen peroxide into water and oxygen:
    \[
    \mathrm{2H_2O_2\longrightarrow 2H_2O+O_2}.
    \]
    \item Glutathione peroxidase reduces hydrogen peroxide using reduced glutathione.
\end{itemize}

\paragraph{Non-enzymatic antioxidants}
\begin{itemize}
    \item Vitamin C.
    \item Vitamin E.
    \item Carotenoids (including beta-carotene).
    \item Glutathione.
    \item Selenium, which is required for the activity of some antioxidant enzymes.
\end{itemize}

\subsection*{7. Prevention and reduction of oxidative stress}
\begin{itemize}
    \item Eat a balanced diet containing fruits and vegetables.
    \item Avoid smoking and second-hand smoke.
    \item Exercise regularly and appropriately.
    \item Limit unnecessary exposure to harmful radiation and toxic chemicals.
    \item Manage chronic conditions with appropriate healthcare.
    \item Avoid taking antioxidant supplements in high doses without professional advice.
\end{itemize}

\subsection*{8. Conclusion}
Reactive oxygen species are produced naturally in the body and have useful roles at controlled levels. Excessive ROS can cause oxidative stress and damage lipids, proteins, DNA, and different body systems. Antioxidant enzymes and dietary antioxidants help control ROS and protect cells.

\end{document}
